\section{双陪集分解与 Hecke 算子的定义}
\label{sec:double-coset}
% ============================================================================

我们现在把上面的几何直觉翻译成精确的代数定义。

\textbf{权 $k$ 斜算子（slash operator）。}
\index{斜算子 (slash operator)}
对矩阵 $\alpha = \begin{pmatrix} a & b \\ c & d \end{pmatrix}$（$\det\alpha > 0$）
和函数 $f\colon \HH \to \C$，定义
\[
  (f \mid_k \alpha)(\tau) = (\det\alpha)^{k/2}(c\tau + d)^{-k} f(\alpha(\tau)).
\]
模形式的变换律 $f(\gamma(\tau)) = (c\tau+d)^k f(\tau)$（$\gamma \in \SL_2(\Z)$）
等价于 $f \mid_k \gamma = f$（对 $\det\gamma = 1$ 时 $(\det\gamma)^{k/2} = 1$）。

斜算子满足 $f \mid_k (\alpha\beta) = (f \mid_k \alpha) \mid_k \beta$。

\begin{definition}[Hecke 算子 $T_n$（全模群情形）]
\label{def:hecke-operator-full}
\index{Hecke 算子 (Hecke operator)}\index{Tn@$T_n$}
对正整数 $n$，令
\[
  \Delta_n = \left\{
    \begin{pmatrix} a & b \\ 0 & d \end{pmatrix}
    \;\middle|\; a, d > 0,\; ad = n,\; 0 \leq b < d
  \right\}.
\]
这是一组矩阵，共 $\sum_{d \mid n} 1$ 个，更精确地共 $\sigma_0(n) = \sum_{d \mid n} 1$ 个。

对 $f \in \Mk_k(\SL_2(\Z))$，定义
\[
  (T_n f)(\tau) = n^{k-1} \sum_{\alpha \in \Delta_n} (f \mid_k \alpha)(\tau)
  = n^{k/2-1} \sum_{\substack{ad=n,\, a,d>0 \\ 0 \leq b < d}}
    d^{-k} f\!\left(\frac{a\tau + b}{d}\right).
\]
\end{definition}

% figures/ch05-hecke-double-coset.tex
% 双陪集分解示意：Γ \ ΓαΓ 的陪集代表元（T_p 的情形）
\begin{figure}[ht]
\centering
\begin{tikzpicture}[>=Stealth, font=\small, scale=1.0]

% 标题
\node[font=\normalsize\bfseries] at (0, 4.2)
  {$T_p$ 的双陪集分解：$\SL_2(\Z) \begin{pmatrix}p&0\\0&1\end{pmatrix} \SL_2(\Z)$};

% 左边：大矩形表示 ΓαΓ（整体）
\draw[thick, blue!50!black, fill=blue!5, rounded corners=6pt]
  (-5.5, -0.3) rectangle (5.5, 3.5);
\node[blue!60!black] at (0, 3.2) {双陪集 $\Gamma\alpha\Gamma$，其中 $\alpha = \begin{pmatrix}p&0\\0&1\end{pmatrix}$};

% p+1 个陪集代表元方块
% 代表元：\begin{pmatrix}1&j\\0&p\end{pmatrix} for j=0,...,p-1 和 \begin{pmatrix}p&0\\0&1\end{pmatrix}

\foreach \j/\xpos/\lab in {
  0/-4.5/{$\begin{pmatrix}1&0\\0&p\end{pmatrix}$},
  1/-2.7/{$\begin{pmatrix}1&1\\0&p\end{pmatrix}$},
  2/-0.9/{$\begin{pmatrix}1&2\\0&p\end{pmatrix}$},
  3/0.9/{$\cdots$},
  4/2.7/{$\begin{pmatrix}1&p{-}1\\0&p\end{pmatrix}$},
  5/4.5/{$\begin{pmatrix}p&0\\0&1\end{pmatrix}$}
}{
  \draw[fill=orange!15, draw=orange!60!black, rounded corners=3pt]
    (\xpos-0.8, 0.2) rectangle (\xpos+0.8, 2.5);
  \node[font=\footnotesize, align=center] at (\xpos, 1.35) {\lab};
}

% 标注个数
\node[font=\footnotesize, text=gray!70!black] at (0, -0.8)
  {共 $p+1$ 个陪集代表元（当 $p$ 为素数时）};

% Hecke 算子作用
\node[font=\footnotesize, align=center, text=purple!70!black] at (0, -1.5) {
  $T_p f(\tau) = p^{k-1} \displaystyle\sum_{j=0}^{p-1} f\!\left(\frac{\tau+j}{p}\right)
  + f(p\tau)$
};

\end{tikzpicture}
\caption{$T_p$ 的双陪集分解。$\SL_2(\Z)$ 在矩阵 $\alpha_p = \begin{pmatrix}p&0\\0&1\end{pmatrix}$
的双陪集 $\SL_2(\Z)\,\alpha_p\,\SL_2(\Z)$ 中，
恰好有 $p+1$ 个右陪集（左陪集代表元），
分别对应矩阵 $\begin{pmatrix}1&j\\0&p\end{pmatrix}$（$j=0,\ldots,p-1$）
和 $\begin{pmatrix}p&0\\0&1\end{pmatrix}$。
权 $k$ 的 Hecke 算子 $T_p$ 就是对这 $p+1$ 个代表元的"平均"。}
\label{fig:hecke-double-coset}
\end{figure}


\textbf{为什么 $\Delta_n$ 是正确的陪集代表元集合？}
$\SL_2(\Z) \backslash M_n(\Z)$（其中 $M_n(\Z) = \{\alpha \in \Mat_2(\Z) \mid \det\alpha = n\}$）
的每个左陪集恰好由 $\Delta_n$ 的某个元素代表。这来自对整数矩阵做初等行变换
（类比于整数矩阵的 Smith 标准型）——对任意行列式为 $n$ 的整数矩阵，
通过左乘 $\SL_2(\Z)$ 中的元素（初等行变换），可以化成上三角形式
$\begin{pmatrix} a & b \\ 0 & d \end{pmatrix}$（$ad = n$，$0 \leq b < d$）。

\begin{example}[$T_2$ 的显式公式（$k = 12$，作用在 $\Delta$）]
\label{ex:T2-on-Delta}
$n = 2$，$\Delta_2$ 包含三个矩阵：
\[
  \begin{pmatrix}1&0\\0&2\end{pmatrix},\quad
  \begin{pmatrix}1&1\\0&2\end{pmatrix},\quad
  \begin{pmatrix}2&0\\0&1\end{pmatrix}.
\]
故（用 $k = 12$）：
\[
  (T_2 f)(\tau)
  = 2^{11}\left[
    2^{-12}f\!\left(\frac{\tau}{2}\right)
    + 2^{-12}f\!\left(\frac{\tau+1}{2}\right)
    + 1^{-12}f(2\tau)
  \right]
  = \frac{1}{2}\left[f\!\left(\frac{\tau}{2}\right)
    + f\!\left(\frac{\tau+1}{2}\right)\right] + 2^{11} f(2\tau).
\]
\end{example}

\begin{example}[$T_p$ 的简洁公式（素数 $p$）]
\label{ex:Tp-formula}
当 $n = p$ 素数时，$ad = p$ 只有 $(a,d) = (1,p)$ 或 $(p,1)$，故
\[
  (T_p f)(\tau) = p^{k-1} \left[
    p^{-k} \sum_{j=0}^{p-1} f\!\left(\frac{\tau+j}{p}\right) + f(p\tau)
  \right]
  = p^{k/2-1}\sum_{j=0}^{p-1} f\!\left(\frac{\tau+j}{p}\right)
    + p^{k/2-1} \cdot p \cdot f(p\tau).
\]
更常见的写法（归一化后）：
\[
  \boxed{(T_p f)(\tau)
  = p^{k-1} f(p\tau) + \frac{1}{p}\sum_{j=0}^{p-1} f\!\left(\frac{\tau+j}{p}\right).}
\]
\end{example}

\begin{definition}[Diamond 算子]
\label{def:diamond-operator}
\index{Diamond 算子 (Diamond operator)}\index{diamond@$\langle d \rangle$}
对同余子群 $\GammaI(N)$，还有另一族算子——\textbf{Diamond 算子}
$\langle d \rangle$（$d \in (\Z/N\Z)^\times$）：

取任意 $\gamma = \begin{pmatrix} a & b \\ c & \delta \end{pmatrix} \in \SL_2(\Z)$，
$\delta \equiv d \pmod N$，定义
\[
  \langle d \rangle f = f \mid_k \gamma.
\]
（选择不同的 $\gamma$ 给出相同的结果，因为它们在 $\GammaI(N)$ 中等价。）

Diamond 算子与 Hecke 算子交换，且 $\langle d \rangle$ 对 $f \in \Mk_k(\GammaI(N), \chi)$
的作用是乘以 $\chi(d)$。
\end{definition}


% ============================================================================
